\chapter[CAPÍTULO VI]{EVALUACIÓN TÉCNICA Y ECONÓMICA DE LA MEJORA / INNOVACIÓN / CREATIVIDAD}

\section{Beneficio técnico y/o económico esperado de la mejora / Innovación / Creatividad}

La implementación del \textbf{Sistema Web Help-Desk Centralizado} en \textbf{J\&P Periféricos S.A.C.} está orientada a resolver el problema principal de desorganización y falta de trazabilidad en la gestión de incidencias y el control del inventario informático multinivel (empresa cliente, sucursal, área y trabajador) para clientes corporativos externos. Desde la ingeniería de software (arquitectura hexagonal en NestJS con TypeScript, Prisma ORM sobre MySQL, MongoDB para historial de chats, Redis Pub/Sub y WebSockets en tiempo real, y \textit{frontend} modular en Vite + React con \textit{shadcn/ui}) y las prácticas de \textbf{ITSM / ITIL v4 e ITAM}, los beneficios esperados se estructuran en dimensiones técnicas y económicas:

\begin{itemize}
    \item \textbf{Centralización y trazabilidad integral de solicitudes (SPOC):} Reemplaza el ingreso informal por llamadas telefónicas y mensajes de WhatsApp sin registro (Causa CA2) mediante una consola web centralizada donde cada incidencia se registra con un código PIN único, máquina de estados auditable y registro de eventos críticos (\texttt{auditoria\_eventos}), asegurando que no se pierdan solicitudes ni se dupliquen atenciones.
    \item \textbf{Agilidad y exactitud en la gestión de activos TI (ITAM):} Sustituye las planillas dispersas y desactualizadas de Excel (Causa CA1) por un expediente digital multinivel con hoja de vida técnica por equipo, trazabilidad de componentes físicos instalados/reemplazados (\texttt{hardware} y \texttt{registro\_hardware}) y bóveda de licencias de software (\texttt{software} y \texttt{software\_equipos}), reduciendo el tiempo de búsqueda previa al diagnóstico.
    \item \textbf{Optimización de visitas técnicas \textit{in-situ} y control de concurrencia:} Incorpora un motor de autoasignación por zona geográfica (\texttt{tecnicos\_zonas}), evaluación de carga laboral, agrupación de hasta tres atenciones en una misma sucursal y franja horaria, y bloqueos temporales (\textit{soft\_locks}), reduciendo los cruces de agenda (Causa CA3), los traslados duplicados y los tiempos muertos en ruta.
    \item \textbf{Estandarización del escalamiento y validación digital de planes y SLA:} Formaliza el flujo de derivación desde soporte remoto hacia soporte \textit{in-situ} con registro del motivo y antecedentes del ticket (Causa CA4), al tiempo que centraliza los planes de soporte (vigencia, servicios incluidos, precios y límite de equipos) para evitar atenciones fuera de cobertura y mejorar el cumplimiento de los Acuerdos de Nivel de Servicio (Causa CA5).
    \item \textbf{Analítica continua y soporte a la decisión gerencial:} Automatiza el cálculo y visualización de indicadores clave de rendimiento (\textbf{MTTA, MTTR, FCR y CSAT}) y métricas operativas en el \textit{dashboard} (Causas CA6 y CA7), eliminando la consolidación manual de reportes a fin de mes y brindando información objetiva para evaluar la calidad del servicio.
    \item \textbf{Retención y confianza de clientes corporativos:} Provee transparencia sobre el estado de las incidencias, el historial de intervenciones y el cumplimiento de tiempos pactados, fortaleciendo la relación comercial de J\&P Periféricos S.A.C. con sus empresas clientes.
\end{itemize}

\section{Relación de beneficio / costo}

La evaluación económica contrasta la \textbf{inversión total de implementación} cuantificada en el Capítulo V (\textbf{S/ 18,254.19} o \textbf{US\$ 5,337.48} al tipo de cambio referencial de S/ 3.42) frente a los \textbf{ahorros operativos} derivados de mitigar las ineficiencias diagnosticadas en el Capítulo III, aplicando el indicador de relación \textbf{Beneficio/Costo (B/C)} y el \textbf{período de recuperación de la inversión (\textit{Payback})}.

\subsection{Identificación de la Pérdida Económica}

A partir del diagnóstico cuantitativo del Capítulo III ---basado en el registro empírico de 5 meses (mayo a septiembre de 2026, con un total de 395 atenciones, una plantilla operativa de 3 técnicos y un costo hora-hombre de \textbf{S/ 25.00})--- se identifican las pérdidas económicas directas asociadas a las tres causas principales priorizadas en el análisis de Pareto, cuyo consolidado mensual y anual se detalla en la Tabla \ref{tab:perdida_economica_cap6}:

\begin{enumerate}
    \item \textbf{Causa CA1 --- Búsqueda manual de datos de equipos en archivos Excel desactualizados:} Se registraron 103 atenciones con demora en 5 meses (26\% promedio), con una paralización estimada de 25 minutos por incidencia. El impacto en los 5 meses asciende a \textbf{S/ 1,072.90}, equivalente a un promedio mensual de \textbf{S/ 214.58} y una proyección anual (12 meses) de \textbf{S/ 2,574.96}.
    \item \textbf{Causa CA2 --- Recepción y comunicación por canales informales (WhatsApp/llamadas) sin registro:} Se registraron 149 solicitudes informales en 5 meses (38\% promedio), con un tiempo improductivo promedio de 30 minutos por solicitud. El impacto en los 5 meses asciende a \textbf{S/ 1,862.50}, equivalente a un promedio mensual de \textbf{S/ 372.50} y una proyección anual (12 meses) de \textbf{S/ 4,470.00}.
    \item \textbf{Causa CA3 --- Descoordinación y reubicación de horarios por cruces en visitas \textit{in-situ}:} Se registraron 81 visitas presenciales afectadas por cruces de agenda y reprogramaciones en 5 meses (20\% promedio), con un tiempo muerto promedio de 20 minutos y sobrecostos operativos de traslado. El impacto en los 5 meses asciende a \textbf{S/ 675.00}, equivalente a un promedio mensual de \textbf{S/ 135.00} y una proyección anual (12 meses) de \textbf{S/ 1,620.00}.
\end{enumerate}

\begin{table}[H]
\centering
\small
\renewcommand{\arraystretch}{1.25}
\caption{Consolidado de Pérdidas Económicas Operativas Identificadas (Situación Actual)}
\label{tab:perdida_economica_cap6}
\begin{tabular}{@{} >{\raggedright\arraybackslash}p{1.2cm} >{\raggedright\arraybackslash}p{5.2cm} >{\raggedright\arraybackslash}p{3.0cm} >{\raggedleft\arraybackslash}p{2.0cm} >{\raggedleft\arraybackslash}p{2.2cm} >{\raggedleft\arraybackslash}p{1.6cm} @{}}
\toprule
\textbf{Código} & \textbf{Causa Raíz Principal} & \textbf{Tipo de Pérdida} & \textbf{Estimación Mensual (S/.)} & \textbf{Proyección Anual (S/.)} & \textbf{Part. (\%)} \\
\midrule
CA1 & Búsqueda manual de datos de equipos en archivos Excel desactualizados & Operativa / Horas-Hombre & 214.58 & 2,574.96 & 29.72\% \\
CA2 & Recepción y comunicación por canales informales (WhatsApp/llamadas) sin registro & Operativa / Horas-Hombre & 372.50 & 4,470.00 & 51.59\% \\
CA3 & Descoordinación y reubicación de horarios por cruces en visitas \textit{in-situ} & Operativa / Logística & 135.00 & 1,620.00 & 18.70\% \\
\midrule
\textbf{Total} & \textbf{Pérdida económica operativa total} & \textbf{Directa (H-H y Traslados)} & \textbf{722.08} & \textbf{8,664.96} & \textbf{100.00\%} \\
\bottomrule
\end{tabular}
\\[0.3em]
{\small \textit{Nota.} Elaboración propia a partir del diagnóstico cuantitativo del Capítulo III (serie mayo--septiembre de 2026 proyectada a 12 meses). Por criterio conservador, no incluye costos indirectos por \textit{downtime} prolongado del cliente ni penalidades.}
\end{table}

\subsection{Identificación de los beneficios del sistema}

Con la puesta en marcha del sistema web centralizado, se estima una \textbf{eficiencia de recuperación del 85.0\%} sobre las pérdidas operativas anuales identificadas. Se asume un \textbf{15.0\% de pérdida residual} debido a factores operativos no eliminables al 100\% en el corto plazo, tales como la curva de aprendizaje de los usuarios finales, eventuales reportes incompletos iniciales o saturación puntual de técnicos de campo ante picos de visitas presenciales:

\begin{equation}
\text{Ahorro Anual Proyectado} = \text{Pérdida Anual} \times \text{Eficiencia de Recuperación}
\end{equation}

\begin{equation}
\text{Ahorro Anual Proyectado} = \text{S/ } 8,664.96 \times 0.85 = \text{S/ } 7,365.22
\end{equation}

\begin{equation}
\text{Pérdida Residual Estimada (15\%)} = \text{S/ } 8,664.96 - \text{S/ } 7,365.22 = \text{S/ } 1,299.74
\end{equation}

La síntesis comparativa entre las pérdidas actuales, el ahorro proyectado, la pérdida residual y la inversión requerida se muestra en la Tabla \ref{tab:resumen_beneficios_cap6}.

\begin{table}[H]
\centering
\small
\renewcommand{\arraystretch}{1.25}
\caption{Resumen de Ahorros Anuales Proyectados e Inversión del Proyecto}
\label{tab:resumen_beneficios_cap6}
\begin{tabular}{@{} >{\raggedright\arraybackslash}p{9.2cm} >{\raggedleft\arraybackslash}p{3.0cm} >{\raggedleft\arraybackslash}p{2.4cm} @{}}
\toprule
\textbf{Concepto Económico / Operativo} & \textbf{Valor (S/.)} & \textbf{Proporción (\%)} \\
\midrule
Pérdida económica anual estimada (situación actual) & 8,664.96 & 100.0\% \\
Ahorro anual proyectado por eficiencia del sistema (85.0\%) & 7,365.22 & 85.0\% \\
Pérdida residual estimada tras la mejora (15.0\%) & 1,299.74 & 15.0\% \\
\midrule
\textbf{Inversión total de implementación (Capítulo V)} & \textbf{18,254.19} & \textbf{---} \\
\bottomrule
\end{tabular}
\\[0.3em]
{\small \textit{Nota.} Elaboración propia a partir de la estimación de ahorro operativo anual y el presupuesto consolidado del Capítulo V.}
\end{table}

\subsection{Relación de beneficio / Costo}

Para determinar el retorno directo por unidad monetaria invertida en el primer año de operación, se calcula la relación \textbf{Beneficio / Costo ($\text{B/C}$)} contrastando el ahorro anual proyectado (\textbf{S/ 7,365.22}) frente al costo total de implementación (\textbf{S/ 18,254.19}):

\begin{equation}
\text{B/C}_{\text{anual}} = \frac{\text{Ahorro Anual}}{\text{Costo de Inversión}} = \frac{7,365.22}{18,254.19} \approx 0.40
\end{equation}

\textbf{Fundamentación e interpretación del resultado:}
\begin{itemize}
    \item En una evaluación estática a \textbf{1 año}, el índice $\text{B/C} = 0.40$ significa que por cada \textbf{S/ 1.00} invertido en la implementación del sistema, se recuperan \textbf{S/ 0.40} en ahorro directo de horas-hombre y logística durante el primer año.
    \item Este resultado obedece a que el proyecto incluye un componente relevante de desarrollo de software a medida (mano de obra de tres desarrolladores en 149 días laborables, que representa el 73.5\% de la inversión total) y equipamiento/soporte, mientras que el ahorro base se calcula de forma conservadora sobre una muestra operativa acotada (395 atenciones en 5 meses).
    \item Si se proyecta el ahorro acumulado en un horizonte operativo de \textbf{3 años} ($\text{S/ } 7,365.22 \times 3 = \text{S/ } 22,095.66$), el beneficio acumulado supera la inversión inicial (ratio acumulado a 3 años de \textbf{1.21}), y a \textbf{5 años} ---alineado con la vida útil de los equipos de cómputo del Capítulo V--- alcanza \textbf{S/ 36,826.10} (ratio acumulado a 5 años de \textbf{2.02}), además de habilitar la comercialización modular externa señalada en el Capítulo I como vía adicional de retorno para J\&P Periféricos S.A.C.
\end{itemize}

\subsection{Relación Costo/beneficio (Tiempo Recuperación)}

El período de recuperación de la inversión (\textit{Payback}) indica el tiempo necesario para que los ahorros operativos acumulados igualen el costo total de implementación (\textbf{S/ 18,254.19}), asumiendo un flujo constante de ahorro anual de \textbf{S/ 7,365.22} (equivalente a \textbf{S/ 613.77} mensuales):

\begin{equation}
\text{Tiempo de Recuperación} = \frac{\text{Costo de Inversión}}{\text{Ahorro Anual}} = \frac{18,254.19}{7,365.22} \approx 2.48\text{ años}
\end{equation}

Convirtiendo la fracción decimal a meses y días:

\begin{equation}
2.48\text{ años} \times 12\text{ meses/año} = 29.76\text{ meses}
\end{equation}

\begin{equation}
29\text{ meses} + (0.76 \times 30\text{ días}) = 29\text{ meses y } 22.8\text{ días} \approx 2\text{ años, } 5\text{ meses y } 23\text{ días}\ (\approx 2.5\text{ años})
\end{equation}

\textbf{Fundamentación e interpretación del resultado:}
\begin{itemize}
    \item Bajo los supuestos conservadores de carga operativa y eficiencia de recuperación del 85\%, J\&P Periféricos S.A.C. recupera la inversión total en aproximadamente \textbf{2 años y 6 meses} (29.76 meses).
    \item A partir de ese punto de equilibrio, los ahorros operativos generados por la trazabilidad de tickets, el expediente ITAM multinivel y la programación sin cruces de visitas \textit{in-situ} representan un beneficio neto recurrente para la empresa, el cual puede acelerarse en la medida en que crezca el volumen de incidencias atendidas y se incorporen clientes corporativos adicionales a la plataforma.
\end{itemize}